\documentclass[../main.tex]{subfiles}

\begin{document}

\chapter{Diskussion}
\vspace{-0.5cm}
\noindent\textit{\small Verfasst von Julian}
\vspace{0.5cm}

Die vorliegende Arbeit verfolgte das Ziel, die Erdrutschanfälligkeit im norditalienischen Alpenraum mithilfe maschineller Lernverfahren zu modellieren und unter verschiedenen Klimaszenarien zu projizieren.
Hierzu wurden zwei Random-Forest-Modelle entwickelt, die im Folgenden hinsichtlich ihrer Stärken, Limitationen und methodischen Implikationen kritisch diskutiert werden.

\section{Bewertung der Modellperformanz}
Beide Modelle erreichen eine ausreichende Modellperformanz, um als Grundlage für operative Entscheidungen im Naturgefahrenmanagement dienen zu können. \\

Dennoch sind beide mit starken Limitationen behaftet, die im Folgenden diskutiert werden.

\subsection{Modell A}
Modell A, entwickelt um die räumliche Verteilung von Erdrutschen in verschiedenen Klimaszenarien abzubilden, erreichte für die Zielklasse Erdrutschereignisse einen Recall von 76\% und einen Precision von 55\%.
Auch wenn dies nicht optimal ist, so ist es dennoch angesichts der limitieren Daten und der Komplexität ein akzeptables Ergebnis.
Gerade bei der Vorhersage über einen großen Zeitraum ist so oder so anzunehmen, dass eine exakte Vorhersage nicht möglich ist.
Daher reicht eine Tendenz aus, vor allem um die verschiedenen Klimaszenarien miteinander vergleichen zu können, was mit Erfolg gelang.

\subsection{Modell B}
Modell B, entwickelt um die tagesaktuelle Erdrutschanfälligkeit abzubilden, erreichte eine akzeptable Modellperformanz mit einem Recall von 0,59\% und einem Precision von 0,62\% für die Zielklasse Erdrutschereignisse.
Das Modell ist wesentlich komplexer als Modell A, da es wesentlich mehr Features verwendet, um die Erdrutschanfälligkeit abzubilden.
Auch wurde and diesem Modell länger und intensiver gearbeitet, als auch getestet.
Verschiedene Features wurden ausprobiert, um die Modellperformanz zu verbessern als auch die Limitationen des Ansatzes zu verstehen. \par

Einer der größten Vorteile zwei Modelle zu verwenden ist, dass der Lernerfolg und damit auch küftige Folgeprojekte durch das Team besser eingeschätzt und bearbeitet werden können.
Modell A ist sehr limitiert, da es an die Daten der Klimaprognose gebunden ist.
Modell B ist da flexibler, da es theoretisch auf beliebige Daten zurückgreifen kann.


\section{Einordnung der Ergebnisse}
Modell A konnte die räumliche Verteilung von Erdrutschen in verschiedenen Klimaszenarien gut abbilden und hat somit die Forschungsfrage beantwortet.
Modell B konnte die tagesaktuelle Erdrutschanfälligkeit gut abbilden und hat somit wichtige Erkenntnisse für die Praxis geliefert. \\
Dementsprechend sind die Ergebnisse der beiden Modelle unterschiedlich einzuordnen, aber dennoch als Erfolg zu werten.
Dennoch besteht großer Verbesserungsbedarf, wenn diese Modelle aktiv in der Praxis eingesetzt werden sollen.
Die Limitationen beider Modelle werden im Folgenden diskutiert. \\

Einer der überraschendsten Ergebnisse war die Umverteilung der Erdrutschanfälligkeit in den Klimaszenarien.
Die Veränderung der Erdrutschanfälligkeit verschiebt sich und erhöt sich, jedoch nicht in den Regionen, in dies von dem Team erwartet wurde.
Dies stellte ein herausforderndes Ergebnis dar, welches in künftigen Arbeiten genauer untersucht werden sollte.


\section{Limitationen und Ausblick}

Zusammenfassend lassen sich folgende zentrale Limitationen identifizieren:
\begin{enumerate}
    \item \textbf{Anzahl der Erdrutschereignisse:} Die Anzahl der Erdrutschereignisse ist für beide Modelle sehr gering, was die Modellperformanz stark einschränkt. Dies ist einer der größten Limitationen der Arbeit und sollte in künftigen Arbeiten genauer untersucht werden. 
    \item \textbf{Fehlende geologische Prädiktoren:} Beide Modelle verzichten auf geotechnische Parameter wie Lithologie, Bodentiefe oder Verwitterungsgrad, die als konditionierende Faktoren für Hanginstabilitäten bekannt sind \autocite{reichenbachReviewStatisticallybasedLandslide2018}.
    \item \textbf{Monolithischer Algorithmus:} Die ausschließliche Verwendung von Random Forest stellt eine Limitation dar, da andere Algorithmen möglicherweise bessere Ergebnisse erzielt hätten.
    \item \textbf{Unvollständige Inventardaten:} Es ist gut möglich, dass die Qualität des IFFI-Erdrutschinventars variiert. Insbesondere kleinere Ereignisse und solche in schwer zugänglichen Gebieten könnten systematisch unterrepräsentiert sein. Dies könnte zu einer Verzerrung der Modellergebnisse führen.
    \item \textbf{Keine sozioökonomische Integration:} Die generierten Karten bilden physische Anfälligkeiten ab, berücksichtigen jedoch keine Expositions- und Vulnerabilitätskomponenten (z.\,B.\,Bevölkerungsdichte, Infrastruktur), die für eine vollständige Risikobeurteilung essenziell wären. Zwar war dies Teil der ursprünglichen Planung und wurde auch Zeitweise eingefügt, jedoch hat es nicht in der finalen Version der Arbeit geschafft aufgrund der hinzugefügten Komplexität.
\end{enumerate}

Trotz dieser Einschränkungen demonstriert die Arbeit, dass ein komplementärer Modellansatz, bestehend aus langfristiger Suszeptibilitätskartierung und tagesaktueller Risikoberechnung, einen Mehrwert für das regionale Naturgefahrenmanagement bieten kann. 
Die identifizierten räumlichen Hotspots auf Gemeindeebene und die prognostizierten Flächenveränderungen unter verschiedenen SSP-Szenarien liefern eine quantitative Grundlage für die Priorisierung von Schutzmaßnahmen und die langfristige Raumplanung im alpinen Raum.
Dennoch ist weitere Forschung und Entwicklung notwendig, um diese Modelle in der Praxis erfolgreich einzusetzen, gerade im Bereich Datenqualität.

\end{document}